% ============================================================
\section{模型的评价、改进与推广}

\subsection{模型的优点}

\begin{enumerate}
  \item \textbf{物理对象和决策时序保持统一。}
  模型围绕负荷、光伏、外网购电和储能建立总量能量平衡，并在各问中连续传递储能状态。统一的状态和能量口径使不同信息条件下的调度方案具有可比性，也便于检查方案是否满足基本物理约束。

  \item \textbf{建模边界与信息条件相匹配。}
  模型根据决策时点能够获得的信息安排预测、计划和调整环节，在保留共同物理结构的同时，分别处理输入不确定性、日内更新和跨日状态传递，避免用同一组信息假定覆盖所有决策阶段。

  \item \textbf{同时检查经济指标与可行性。}
  各方案以购电及相关附加费用作为主要经济指标，同时复核能量平衡、储能状态、容量和功率边界。经济指标与约束校验同时保留，使评价既能反映运行成本，也能排除不可执行的调度结果。

  \item \textbf{显式表示储能的跨日作用。}
  问题四将日末储能保留量对后续运行的影响表示为终端价值，并在日前计划和日内滚动调整中使用，使储能决策不再局限于单日费用最小化。
\end{enumerate}

\subsection{模型的不足}

\begin{enumerate}
  \item \textbf{聚合模型未描述配电网空间约束。}
  当前模型将小区视为单一聚合节点，未考虑节点电压、线路容量、潮流分布和线路损耗；当局部网络约束成为主要瓶颈时，模型可能低估实际运行限制。

  \item \textbf{储能模型未计入全生命周期因素。}
  模型按固定性能参数描述储能，未考虑温度、老化、自放电、健康状态和循环退化成本。因此，费用结果主要反映运行阶段的购电成本，不等同于储能设备的全生命周期经济收益。

  \item \textbf{预测模型依赖历史规律，面对突发变化的能力有限。}
  问题二、三使用历史信息构造负荷和光伏预测场景，问题四的电价采用前一周同日同刻预测。面对节假日、天气突变或价格机制变化时，历史规律可能不足以代表当前情形，且本文未建立完整的电价概率预测。

  \item \textbf{终端价值和实际执行仍是近似处理。}
  问题四的终端价值由有限预测视野、有限储能取点和小时聚合估计。问题三以及问题四的滚动方案会在规定决策时刻重新优化计划，但这些过程仍是在历史数据上的离线回测，执行阶段按照代码设定的实际供需替换和富余先充、缺口先放规则运行，尚未接入真实能源管理系统进行在线数据采集、控制下发和闭环运行。
\end{enumerate}

\subsection{模型的改进与推广}

\begin{enumerate}
  \item \textbf{引入配电网网络约束。}
  对于园区或多节点微电网，可在现有能量平衡模型上增加节点、线路、潮流和电压约束，将聚合模型扩展为配电网最优潮流或线性化潮流模型。原有储能和费用目标可以保留，但网络参数需要根据实际拓扑重新估计。

  \item \textbf{加入储能退化和健康状态模型。}
  可在目标函数中加入循环退化成本，并根据充放电深度、温度和荷电状态建立分段线性退化模型；若有长期运行数据，还可进一步引入 SOH 状态变量。这样得到的目标将从单纯运行购电成本扩展为运行成本与寿命成本的综合权衡。

  \item \textbf{完善不确定性处理和终端价值估计。}
  可将负荷、光伏和电价的预测误差统一构造为带相关性的联合场景，并针对节假日、天气突变和价格异常增加校准机制；同时采用更细的时间粒度和更充分的储能状态取点估计终端价值，使其与滚动控制更紧密衔接。

  \item \textbf{拓展应用场景。}
  在补充网络拓扑、设备性能和运行数据后，本文的总量平衡与滚动调度思路可推广到园区微电网、工商业储能和光储充系统，也可扩展到多储能协调及需求响应问题。迁移时需要重新确定负荷、价格、储能和网络参数，并依据实际通信与控制条件重新检验可执行性。
\end{enumerate}
